% research-program.tex — Section 7: The Research Program

\section{The Research Program}
\label{sec:research-program}

\subsection{What is explained, what is conditional, what is contingent}
\label{subsec:status}

\begin{table}[ht]
\centering
\caption{Status of physical content within the program.}
\label{tab:status}
\small
\begin{tabular}{@{}p{3.2cm}p{1.6cm}p{2.5cm}@{}}
\toprule
Content & Status & Source \\
\midrule
Complex QM & \textbf{Forced$^*$} & Paper~5 (A1--A4) \\
Born rule & \textbf{Forced} & Gleason ($n \geq 3$) \\
Unitary dynamics & \textbf{Forced} & Skolem-Noether \\
Area-law entanglement & \textbf{Forced$^*$} & Paper~6 (UV layer) \\
Einstein's equations & Conditional & Paper~6 (4 std gaps) \\
SM gauge group & Conditional & Paper~7 (Gap~C) \\
Chiral representation & Conditional & Paper~7 (Gap~C) \\
3 generations & Speculative & Peirce tower \\
Mass hierarchy & Speculative & Peirce depth \\
Specific couplings & Contingent & Ring in structure space \\
3+1 dimensions & Plausible & Huygens + stability \\
\bottomrule
\end{tabular}
\end{table}

The program's intellectual honesty requires distinguishing these
categories sharply.  ``Forced'' means the result follows by theorem
from the premises.  ``Conditional'' means the result holds subject to
identified gaps that are registered explicitly.  ``Speculative'' means
a structural argument exists but has not been formalized.
``Contingent'' means the result may not be uniquely determined by
self-modeling.

Specific coupling constants, masses, and mixing angles may belong to
the contingent category: the self-modeling condition constrains the
algebra ($A_F$, forced) but not the dynamics ($D_F$, where parameters
live).  The viable parameter values likely form a manifold --- a
``ring around the basin'' --- not a unique point.  The experiential
measure $\rho$ (Remark~\ref{rem:rho-h3o}) selects a position on this
ring, but the selection may
have the same status as ``why this brain?'' rather than ``why
gravity?''  3+1 dimensions and the SM gauge group are likely universal
across the ring (forced by algebraic structure).  The fine structure
constant may vary.

\subsection{Falsifiable predictions}
\label{subsec:predictions}

\begin{enumerate}
\item \textbf{No coupling unification (prediction).}  The SM gauge
  group arises from an intersection of stabilizers in $\Ffour$, not
  from a single GUT group breaking.  The framework predicts that
  running couplings do \emph{not} unify at a single scale.

\item \textbf{No Planck-scale modifications to QM (prediction).}
  Quantum mechanics is forced by the algebra of self-modeling
  (Steps~3--6), not by spacetime structure.  There is no reason for QM
  to break down at the Planck scale.  Searches for Planck-scale
  decoherence should find null results.

\item \textbf{Proton stability (consistency check).}  No $\SU{5}$ or
  $\mathrm{SO}(10)$ GUT group appears in the framework.  Proton decay
  via GUT-mediated channels is not predicted.  This is consistent with
  current bounds ($\tau_p > 10^{34}$ years).  Current data on coupling
  non-unification in the non-SUSY SM is also consistent
  with the intersection-of-stabilizers mechanism.

\item \textbf{Cubic vs.\ quadratic consciousness measure (prediction).}
  $\rhoJ$ (Claim~\ref{claim:rhoJ}) depends on $\det(X)$, a cubic
  invariant.  IIT's $\Phi$~\citep{tononi2008} is quadratic.
  Selectively disrupting one of three Peirce sectors should produce a
  sharp (not gradual) loss of self-modeling quality.  Thalamocortical
  disruption data from anesthesia research can distinguish these
  functional forms.
\end{enumerate}

\subsection{Open problems}
\label{subsec:open-problems}

\begin{enumerate}
\item \textbf{Gap~C (complexification).}  The most pressing open
  problem.  A theorem of the form ``$C^*$-measurement maps on a real
  Jordan module necessarily factor through $\otimes_{\mathbb{R}}
  \mathbb{C}$'' would close it.  Failing that, the evolutionary
  argument stands but the SM derivation remains conditional.

\item \textbf{Spectral action on $\hthree$.}
  Farnsworth~\citep{farnsworth2015} has computed spectral actions on
  non-associative geometries, producing Einstein gravity coupled to
  $G_2$ gauge fields from octonionic spectral triples.  Extending this
  to $\hthree$ would unify the GR route (Paper~6) with the SM route
  (Paper~7) under a single spectral action.  The computation is open.

\item \textbf{Vanchurin dynamics.}  Vanchurin~\citep{vanchurin2020,
  vanchurin2024} showed that neural network learning dynamics reproduce
  aspects of quantum mechanics and general relativity.  The present
  program contributes the learning objective: self-modeling.  The loss
  function is $-I(B;\,M)$.  Characterizing the fixed point of
  Vanchurin's dynamics with this loss would address Level~6 questions
  (specific parameter values).

\item \textbf{3+1 dimensions.}  3+1 is the unique dimensionality
  where clean signal propagation (Huygens' principle), stable atoms,
  stable orbits, and knot existence all hold simultaneously.  The
  Huygens argument is specifically about self-modeling: a self-modeler
  receiving echoes of old signals double-counts information it already
  processed.  Formalizing this as a channel capacity argument is open.

\item \textbf{Infinite-dimensional extension.}  The full chain is
  proved for finite-dimensional spectral order-unit spaces.  Quantum
  field theory is not finite-dimensional.  The Bekenstein bound
  suggests that self-modelers in finite regions are necessarily
  finite-dimensional, but connecting this to the formalism is open
  work.

\item \textbf{Lean formalization.}  Paper~5's chain has been
  scaffolded in Lean~4.  Completing the formalization (removing all
  \texttt{sorry} placeholders) would upgrade the chain from
  ``verified by human review and adversarial testing'' to ``machine
  verified.''
\end{enumerate}

\subsection{Closing}
\label{subsec:closing}

The Copernican program has a simple logic: find a distinction that is
doing no work, remove it, and see what follows.  Each removal
simplified the theory and expanded its explanatory scope.  We have
argued that the distinction between ``mathematical'' and ``physical''
is the last such distinction, and that removing it forces the physics
observers experience --- from the complex field to the Standard Model
gauge group --- through the algebraic consequences of self-modeling.

The chain is not complete.  Gap~C is open.  The spectral action on
$\hthree$ is uncomputed.  The Vanchurin dynamics have not been solved
with self-modeling as the loss function.  But the chain is long enough,
and the gaps are honest enough, to constitute a research program worth
pursuing.  The program's distinguishing feature is not that it has
fewer gaps than its competitors, but that it has \emph{identified}
gaps where others have hidden assumptions, and that its foundational
premise --- no privileged instantiation --- is not an axiom added to
physics but the last axiom removed from it.
